\documentclass[11pt]{article}
\usepackage[T1]{fontenc}
\usepackage[utf8]{inputenc}
\usepackage{lmodern,amsmath,amssymb,amsthm,mathtools,booktabs,longtable,array}
\usepackage[margin=1in]{geometry}
\usepackage[hyphens]{url}
\usepackage[colorlinks=true,linkcolor=blue,citecolor=blue,urlcolor=blue]{hyperref}
\usepackage{microtype}
\hypersetup{pdftitle={Joint convexification: exact support, complete tolerance separation, and original-model cuts},pdfauthor={}}
\newtheorem{theorem}{Theorem}[section]
\newtheorem{lemma}[theorem]{Lemma}
\newtheorem{proposition}[theorem]{Proposition}
\newtheorem{corollary}[theorem]{Corollary}
\theoremstyle{definition}
\newtheorem{definition}[theorem]{Definition}
\theoremstyle{remark}
\newtheorem{remark}[theorem]{Remark}
\newcommand{\R}{\mathbb R}
\newcommand{\Q}{\mathbb Q}
\newcommand{\E}{\mathbb E}
\newcommand{\conv}{\operatorname{conv}}
\newcommand{\supp}{\operatorname{supp}}
\newcommand{\argmin}{\operatorname{argmin}}
\newcommand{\norm}[1]{\lVert#1\rVert}
\newcommand{\pathref}[2]{\href{run:#1}{\texttt{\detokenize{#2}}}}
\title{Joint convexification: exact support, complete tolerance separation, and original-model cuts}
\author{}
\date{October 2026}
\begin{document}
\maketitle
\begin{abstract}
Joint convexification uses the dependence among nonlinear expressions that
share variables and domain constraints. This report develops the support,
separation, representation, and verification steps needed to turn that
principle into a specified solver component. It gives exact rational support
for quadratic objectives on bounded rational polytopes, retains a polynomial
support algorithm for constrained quadratic stars, and supplies a finite
positive-tolerance separation procedure with explicit costs and failure
statuses. An original-variable aggregation formula transfers joint support
bounds to linear cuts without introducing nonlinear feature equalities.
The implementation checks original expression domains, model binding, and
the final floating-point row. In a new prospective 30-model application
sample, native baseline and both cut modes solve the same 25 models within
the 30-second budget. Cut modes use more summed time; the result does not
support default activation. Native SCIP remains the default. A separate
75-job matched validation checks repairs for four retained diagnostic
worker errors and discovery-budget violations, with no extra solves.
All 165 recorded cuts across both campaigns pass replay; the four original
missing logs remain excluded. Exact support, complete
tolerance separation, numerical solver
performance, and complete certification of a solve are distinct claims.
The report proves and implements the first two on their stated classes and
evaluates the third; it does not claim the fourth or publication priority
for established convexification, duality, or quadratic-programming methods.
\end{abstract}
\tableofcontents

\section{The problem and the completion standard}\label{sec:scope}

Several functions of the same variables can be convexified together by
minimizing their weighted sum over the common domain. This familiar
observation leaves several practical questions. Which domains and functions
have a complete support algorithm? Does a direction search actually separate
a point, or merely fail to find a cut? Can the cut be applied while retaining
the native nonlinear rows? Does the exact inequality survive floating-point
assembly? Does the resulting solver solve more problems, or spend more time
on work that its native relaxations already accomplish?

The previous repository continuation supplied certified support cuts,
exact coupled-pair and constrained-star oracles, automatic graph-block
discovery, and a short controlled experiment. It left native SCIP as the
recommended default: among 24 selected application models, 20 passed the
common importer, baseline and reformulation control solved 19, and each
cut mode solved 18. It also exposed three domain-admission refusals and
one unsupported variable-exponent model. These findings motivate the
present work; they are not silently replaced by its new experiment.
The original report and raw results remain in
\pathref{../../research-20261002-convexification/document/main.pdf}{research-20261002-convexification/document/main.pdf}
and
\pathref{../../research-20261002-convexification/experiments/campaign-v1/records.jsonl}{the first campaign records}.

Completion has a concrete meaning here. The supported algorithms must have
explicit inputs and statuses; proofs must cover their complete admitted
domains, including degeneracy; returned certificates must bind the actual
model and row; a complete separation mode must distinguish a checked
positive-tolerance result from an exhausted heuristic; and empirical
claims must match a frozen, reproducible comparison. A theorem that an
arbitrary nonlinear problem can be solved efficiently is not part of that
standard. Section~\ref{sec:frontier} proves an obstruction to that stronger
interpretation rather than treating it as an omitted implementation step.

The mathematical chain is self-contained. Sections~\ref{sec:support} and
\ref{sec:numerics} restate the inherited support and certification foundation.
Sections~\ref{sec:aggregation}--\ref{sec:separation} develop the current
extensions and retain the constrained-star and overlap results needed to
understand their scope. Section~\ref{sec:implementation} binds those results
to the actual implementation. The
\pathref{COVERAGE.md}{claim and evidence map} distinguishes proved results,
implemented APIs, internal review, numerical checks, and measured outcomes.
Internal review is not formal verification or journal peer review.

% Mathematical foundations adapted from the 2026-10-02 report.
\section{Joint support cuts}\label{sec:support}

Let $D\subset\R^n$ be nonempty and compact, and let
$F=(f_1,\ldots,f_k):D\to\R^k$ be continuous. Write
\[
G_D=\{(x,F(x)):x\in D\},\qquad H_D=\conv G_D.
\]
The lifted variables $w_j$ represent the \emph{same} functions $f_j$ of
the \emph{same} variables $x$. For a direction $d=(a,b)\in\R^{n+k}$ define
\begin{equation}\label{eq:support}
h_D(a,b)=\min_{x\in D}\{a^\top x+b^\top F(x)\}.
\end{equation}
We use lower support values, so every cut is written with a lower bound.
If integer variables occur in $D$, support over a continuous superset
remains valid for the mixed-integer graph, but need not describe its hull.

\begin{proposition}[Support and complete hull description]\label{prop:support}
If $\underline h\le h_D(a,b)$, then
\begin{equation}\label{eq:cut}
a^\top x+b^\top w\ge\underline h
\end{equation}
is valid for $H_D$. Moreover,
\[
H_D=\bigcap_{(a,b)\in\R^{n+k}}
\{(x,w):a^\top x+b^\top w\ge h_D(a,b)\}.
\]
\end{proposition}
\begin{proof}
The defining minimum makes the inequality valid at each point of $G_D$;
linearity preserves it under convex combinations. The graph is compact
by continuity, and its convex hull is compact in finite dimension. A point
outside that hull is strictly separated by a linear functional. Orient
that functional as a lower support inequality to obtain a member of the
displayed intersection that excludes the point.
\end{proof}

This is classical convex hull duality; simultaneous graph separation
from linear combinations is developed explicitly in
\cite{LiersEtAl2021,Tawarmalani2010}. The computational problem is to
evaluate or bound~\eqref{eq:support}, and to find a useful direction
without spending more than its benefit. A numerical direction generator
is permitted to be inaccurate: its output only chooses $(a,b)$.
Soundness depends on a bound for that \emph{final} direction, not on the
correctness or optimality of the direction search. The same observation
allows random samples, LP direction search, and a native sampling kernel
outside the cut-validity calculation.

For $n=1$ and $k=1$, optimizing over the coefficient of $x$ recovers the
biconjugate description of a scalar convex envelope. For $k>1$, joint
directions retain dependencies between outputs that intersected scalar
envelopes can omit. A finite collection of valid joint cuts generally
gives an outer approximation; absence of a violated generated cut does
not establish membership in $H_D$.

\begin{proposition}[Final floating-point row]\label{prop:round}
Suppose the exact values of the coefficients submitted to a numerical
solver are $\widehat a\in\Q^n$ and $\widehat b\in\Q^k$ (in particular,
finite binary floating-point numbers interpreted exactly). If a checker
establishes
\[
L\le\min_{x\in D}
\{\widehat a^\top x+\widehat b^\top F(x)\}
\]
and the submitted right-hand side $\widehat L$ satisfies
$\widehat L\le L$ exactly, the submitted mathematical row is valid on
$H_D$.
\end{proposition}
\begin{proof}
Apply Proposition~\ref{prop:support} to the submitted coefficients and
then weaken the right-hand side from $L$ to $\widehat L$.
\end{proof}

Rounding an already proved rational row without checking its resulting
coefficients would not satisfy this contract. One alternative, when a
trusted box $\ell\le z\le u$ is available for all lifted coordinates,
is to replace the exact row $c^\top z\ge L$ by
\[
\widehat c^\top z\ge
L+\sum_i\min\{(\widehat c_i-c_i)\ell_i,
                 (\widehat c_i-c_i)u_i\},
\]
and round the right-hand side downward. The formula follows by bounding
each coordinate of $(\widehat c-c)^\top z$. Our principal contract is
the direct check for the final direction in
Proposition~\ref{prop:round}; it avoids needing auxiliary-variable bounds
for this correction.

\section{Finite certificates on the whole domain}\label{sec:numerics}

\subsection{Exact Bernstein bounds for rational polynomials}

Consider a polynomial $p$ on a rational box
$B=\prod_{i=1}^n[\ell_i,u_i]$. Substitute
$x_i=\ell_i+(u_i-\ell_i)t_i$ and eliminate fixed coordinates. In the
remaining coordinates, write
\[
p(x(t))=\sum_{\beta\le d}a_\beta t^\beta
       =\sum_{\alpha\le d}b_\alpha B_\alpha^d(t),\qquad
B_\alpha^d(t)=\prod_i
\binom{d_i}{\alpha_i}t_i^{\alpha_i}(1-t_i)^{d_i-\alpha_i}.
\]
Here $d$ is a componentwise degree bound and all coefficients are exact
rationals. The conversion is
\begin{equation}\label{eq:bernstein-conversion}
b_\alpha=\sum_{\beta\le\alpha}a_\beta
\prod_i\frac{\binom{\alpha_i}{\beta_i}}{\binom{d_i}{\beta_i}}.
\end{equation}

\begin{lemma}[Bernstein enclosure]\label{lem:bernstein}
On $B$, $\min_\alpha b_\alpha\le p\le\max_\alpha b_\alpha$.
\end{lemma}
\begin{proof}
For $0\le t_i\le1$, every $B_\alpha^d(t)$ is nonnegative and
$\sum_\alpha B_\alpha^d(t)=1$, by the binomial theorem in each
coordinate. Thus $p(x(t))$ is a convex combination of the coefficients.
Identity~\eqref{eq:bernstein-conversion} follows by expressing each
monomial $t^\beta$ in the tensor-product Bernstein basis; the univariate
identity has coefficients $\binom{\alpha}{\beta}/\binom d\beta$.
\end{proof}

The enclosure is classical
\cite{GarloffSmith2008,GarloffJanssonSmith2003}. The implementation contribution is to bind
the exact polynomial, its original box, and the checked final row to
the recorded partition and lower bounds. Adaptive subdivision can
improve a weak enclosure without changing the validity argument.

\begin{proposition}[Partition certificate]\label{prop:partition}
Let $B_1,\ldots,B_N$ be finitely many closed boxes whose union is $B$,
and suppose $L_j\le\inf_{B_j}p$ is verified for each $j$. Then
$L=\min_jL_j\le\inf_Bp$. The same conclusion holds for a finite interval
partition and any expression with a verified interval enclosure on every
cell.
\end{proposition}
\begin{proof}
Every $x\in B$ belongs to a checked cell $B_j$, where
$p(x)\ge L_j\ge\min_iL_i$.
\end{proof}

The coverage premise is essential. Checking all supplied cells while
accepting an omitted end interval would not prove the claimed bound.
The certificate checker must start from the trusted original domain and
verify the entire subdivision, including endpoints. A cache key also
needs the expression, direction, and domain: reusing a valid bound from a
smaller domain on a larger domain is unsound.

\subsection{Supported elementary expressions and explicit failure}

For a supported elementary expression, outward ball or interval
arithmetic supplies an enclosure for the support objective on each
interval. Rational constants and final binary coefficients are embedded
exactly. Expression evaluation must retain domain requirements: a
logarithm needs a positive argument; division needs a denominator
excluding zero; powers require their stated real domain. Algebraic
simplification is not permission to remove the original expression's
domain restrictions.

If a pole, invalid domain, nonfinite enclosure, unsupported operation, or
resource limit prevents a complete bound, the procedure returns that
status. An exhausted direction search means only that no checked useful
cut was obtained within its budget. It does not prove that a point lies
in a hull, that the original constraints hold, or that the domain is
empty. Even a successfully checked cut need not be violated enough to
justify insertion into the LP.

The elementary support calculation can also use a curvature correction.
If the scalar direction $p$ is twice continuously differentiable on
$[a,b]$ and a verified bound gives $p''\le M$, then
\begin{equation}\label{eq:chord-correction}
\min_{[a,b]}p\ge
\min\{p(a),p(b)\}-\max\{0,M\}\,(b-a)^2/8.
\end{equation}
Indeed, $p(t)-Mt^2/2$ is concave, so its chord inequality rearranges to
$p(t)\ge\operatorname{chord}_p(t)-M(t-a)(b-t)/2$.
Use $(t-a)(b-t)\le(b-a)^2/4$ when $M\ge0$; for $M<0$, the
chord itself is a lower bound. Verified lower endpoint values can
replace $p(a)$ and $p(b)$. The implementation first checks all feature
domains, including zero-weight features, then uses the maximum of this
bound and available natural enclosures. An unavailable derivative bound
or singular derivative disables this optional strengthening; it does
not license removing a domain endpoint.

The supported grammar, actual certificate schema, and solver binding
are stated with their implementation in Section~\ref{sec:implementation}.
Neither Python's rational arithmetic nor the interval library is
formally verified here. Replay reduces the trusted computational path
and catches mismatched artifacts; it is not a proof of the underlying
software stack.


\section{Joint cuts in the original variables}\label{sec:aggregation}

Let the original model contain the inequality sides
\begin{equation}\label{eq:sides}
g_i(x)+\ell_i(v)\le b_i,\qquad i=1,\ldots,m,
\end{equation}
where $x$ is a selected subset of original variables, $v$ contains model
variables, and each $\ell_i$ is linear. A variable may occur both in $x$
and in an affine remainder. All constants are included in $b_i$.
The block domain $D$ contains the block coordinates of every feasible
point to which a proposed cut will apply.

\begin{proposition}[Original-variable support cut]\label{prop:aggregate}
Choose $a\in\R^{\dim x}$ and $\lambda\in\R^m_+$. If
\[
\beta\le\inf_{u\in D}\left\{a^\top u+
                         \sum_i\lambda_i g_i(u)\right\},
\]
then the original-variable linear inequality
\begin{equation}\label{eq:aggregate}
a^\top x-\sum_i\lambda_i\ell_i(v)
\ge\beta-\sum_i\lambda_i b_i
\end{equation}
is valid for the model.
\end{proposition}
\begin{proof}
At a feasible point, nonnegativity of the multipliers gives
$\sum_i\lambda_i g_i(x)\le\sum_i\lambda_i(b_i-\ell_i(v))$.
Substitute this upper bound into the support inequality and rearrange.
\end{proof}

This is weak Lagrangian duality. It needs neither convex functions nor
optimal multipliers. The functions are minimized jointly before elimination;
simply combining already linear valid inequalities cannot strengthen their
intersection. Equality sides can use unrestricted multipliers, or two
opposite inequality sides with nonnegative multipliers. A minimization
objective $g_0(x)+\ell_0(v)+c_0$ supplies an epigraph side
$g_0(x)+\ell_0(v)-t\le-c_0$. A maximization objective uses the
hypograph side $-g_0(x)-\ell_0(v)+t\le c_0$. The objective
transformation, constant, and sense must be identical in every compared
mode. These rows do not require a separate equality variable for each
selected nonlinear expression.

\begin{theorem}[Closure of the direct-row interface]\label{thm:closure}
Suppose $D$ is nonempty and compact and $g:D\to\R^m$ is continuous. Put
\[
K=\conv\{(x,g(x)):x\in D\},\qquad
P=\{(x,z):\exists y,\ (x,y)\in K,\ y+Lz\le b\}.
\]
Then $P$ is exactly the intersection, over all $a\in\R^{\dim x}$ and
$\lambda\in\R^m_+$, of
\begin{equation}\label{eq:closure}
a^\top x-\lambda^\top Lz\ge
\min_{u\in D}\{a^\top u+\lambda^\top g(u)\}-\lambda^\top b.
\end{equation}
If continuous equality features $h:D\to\R^r$ are included in the
joint graph, their support coefficients are free and their target values
are substituted without adding cone coordinates.
\end{theorem}
\begin{proof}
The set $U=K+(\{0\}\times\R^m_+)$ is closed and convex. To see
closedness, take a convergent sequence of sums, extract a convergent
subsequence of the compact $K$ summands, and subtract; the cone summands
then converge inside their closed cone. Membership in $P$ is equivalent
to $(x,b-Lz)\in U$.

A lower support functional of $U$ with finite constant must have
nonnegative coefficients on the last $m$ coordinates, because a negative
coefficient has infimum $-\infty$ along the corresponding cone ray.
For the remaining directions its support value is
$\min_{u\in D}(a^\top u+\lambda^\top g(u))$. Substituting
$(x,b-Lz)$ gives~\eqref{eq:closure}. Closed convex separation excludes
every point outside $U$ by such a finite support inequality, proving the
converse. For continuous equality features the joint graph
$\conv\{(x,g(x),h(x)):x\in D\}$ is also compact. The same argument
adds no cone in their coordinates, so their coefficients are free.
\end{proof}

This describes the projection of the specified graph relaxation, not
necessarily the convex hull of original feasible points. For example,
on $[0,1]$ the equality $x^2=1/4$ has feasible set $\{1/2\}$.
The mixture with mass $3/4$ at zero and $1/4$ at one instead has
$\E x=\E x^2=1/4$. The projected graph relaxation thus contains
$x=1/4$, and every direct support cut accepts it. Including $x^2=1/4$
inside the support domain changes the oracle problem and removes this
example. Convexifying a graph and imposing constraints in expectation do
not generally commute. Ordinary domain propagation can also resolve this
particular equality; that does not change the comparison of cut families
over the unchanged box.

\begin{proposition}[Safe export after elimination]\label{prop:elimination-round}
Let the exactly eliminated row be $c^\top v\ge r$ and let
$\widehat c$ denote the exact rational values of the coefficients
actually exported in binary64. With $\Delta_j=\widehat c_j-c_j$, suppose
valid global bounds imply a finite value
\[
E=\sum_j\inf_{L_j\le s\le U_j}\Delta_j s.
\]
Then every exported lower side $\widehat r\le r+E$ gives a valid row
$\widehat c^\top v\ge\widehat r$.
\end{proposition}
\begin{proof}
$\widehat c^\top v=c^\top v+\sum_j\Delta_jv_j\ge r+E$.
\end{proof}

A positive $\Delta_j$ needs a finite lower bound, a negative one a
finite upper bound, and a zero one contributes zero even on an unbounded
coordinate. This covers an unbounded objective epigraph when its exported
coefficient is exact. The correction must follow coefficient merging,
normalization, and any zero dropping. The lower side is rounded downward
and checked by exact rational comparison. An unavailable finite correction,
overflow, or lost useful separation causes refusal; a fixed safety shift
cannot replace the calculation. The source sides, their orientations,
multipliers, domain provenance, support bound, exact eliminated row, and
actual exported row are separate parts of the certificate contract.

The derivation and its checked implementation are detailed in
\pathref{../theory/row-aggregation.md}{theory/row-aggregation.md}.
The formulation and safe-rounding principles have direct prior literature;
Section~\ref{sec:literature} records the sources and contribution boundary.

\section{Exact quadratic support on a bounded polytope}\label{sec:polytope}

Let $P=\{x\in\R^d:Ax\le b\}$ be a bounded rational polytope and
write a rational quadratic as
\[
q(x)=\tfrac12x^\top Hx+c^\top x+c_0,\qquad H=H^\top.
\]
The support of any vector of quadratic features has this form. The matrix
$H$ may be indefinite or singular. Finite supplied variable bounds make
boundedness part of the checked input; equalities are represented by their
two inequality sides.

For each subset $S$ of at most $d$ rows, form the stationary-face system
\begin{equation}\label{eq:face-system}
\begin{pmatrix}H&A_S^\top\\A_S&0\end{pmatrix}
\begin{pmatrix}x\\\lambda\end{pmatrix}
=\begin{pmatrix}-c\\b_S\end{pmatrix}.
\end{equation}
Retain its unique solution if the matrix is nonsingular and $Ax\le b$
holds exactly. Multipliers are unrestricted: this is stationarity on an
affine face, and sign filtering is unnecessary.

\begin{theorem}[Complete finite quadratic support]\label{thm:polytope}
If $P$ is nonempty, at least one retained candidate minimizes $q$ globally.
The least candidate value is therefore the exact rational support value,
attained at a rational point. If no candidate is retained after the complete
enumeration, $P$ is empty. These statements include lower-dimensional
polytopes, implicit equalities, redundant rows, and constant objectives.
\end{theorem}
\begin{proof}
Compactness gives a global minimizer. Among global minimizers choose $x_*$
whose minimal containing face $F$ has the smallest dimension. It lies in
the relative interior of $F$. Let $T$ be the tangent space of the affine
hull of $F$. Relative first- and second-order conditions imply
$v^\top(Hx_*+c)=0$ and $v^\top Hv\ge0$ for every $v\in T$.

If the Hessian restricted to $T$ were singular on a positive-dimensional
face, there would be a nonzero $v\in T$ with $v^\top Hv=0$.
The quadratic identity and stationarity give
$q(x_*+tv)=q(x_*)$ for every real $t$. The line meets the bounded
face in a closed interval with $x_*$ in its relative interior. At an
endpoint it reaches a proper face with the same minimum, contradicting
the choice of $F$. Thus $H$ is positive definite on $T$ unless $T=\{0\}$.

Choose an independent basis $S$ of the active rows at $x_*$. It has at
most $d$ rows and $\ker A_S=T$. In a homogeneous solution $(v,\mu)$
of~\eqref{eq:face-system}, $A_Sv=0$ and
$Hv+A_S^\top\mu=0$. Multiplication by $v^\top$ gives
$v^\top Hv=0$, hence $v=0$ by positive definiteness on $T$.
Independence of the rows then gives $\mu=0$. The matrix is nonsingular,
including the vertex case $T=\{0\}$. Relative stationarity places
$Hx_*+c$ in the row span of $A_S$, so its unique solution includes
$x_*$. The complete enumeration therefore retains a minimizer.
All systems and evaluations are rational. Conversely, every retained
candidate is feasible, so an empty polytope retains none.
\end{proof}

If there are $m$ rows including variable bounds, there are at most
$N=\sum_{k=0}^d\binom{m}{k}$ subsets. Direct exact elimination on
systems of order at most $2d$ and feasibility testing against all rows
uses $O(N(d^3+md))$ rational arithmetic operations. Rational determinant
bounds control the encoding lengths of candidates and objective values.
For each fixed $d$ the bit complexity is polynomial in the input length;
when $d$ is part of the input, the subset count can be exponential.
This is a classical face-enumeration algorithm
\cite[Section 2.9]{Murty1997}, not a new tractability result for
unrestricted quadratic programming.

The implementation
\pathref{../theory/quadratic_polytope.py}{theory/quadratic_polytope.py}
checks its enumeration budget before claiming completion. A resource
refusal is not an optimum or an emptiness certificate. Its certificate
binds bounds, rows, and coefficients and records the complete candidate
calculation. Replay reconstructs that calculation from trusted inputs;
it shares exact arithmetic routines with generation. Independent tests
and mathematical review supply different checks, but not formal
verification of the Python interpreter or rational-arithmetic library.

The previous two-variable polygon algorithm is a more specialized version:
it checks vertices, stationary edge points, and a positive-definite interior
stationary point in $O(m^3)$ arithmetic operations. The new oracle permits
arbitrary affine coupling in every admitted small dimension. It can thus
merge overlapping quadratic blocks with leaf--leaf rows or products when
the resulting block fits its budget. The next section retains a different
advantage: constrained stars remain polynomial in their number of leaves
without fixing that dimension.

% Adapted from the 2026-10-02 report; inherited result and implementation.
\section{A larger exact class: quadratic stars with affine rows}\label{sec:star}

The obstruction in Section~\ref{sec:overlap} motivates support directions
that span adjacent pairs. A tractable extension is a center variable $y$
with leaves $x_1,\ldots,x_k$. The objective is
\begin{equation}\label{eq:star}
q(y,x)=a_0+a_1y+a_2y^2+
\sum_{i=1}^k\{d_ix_i^2+(e_iy+f_i)x_i\}.
\end{equation}
All objective coefficients, affine-row coefficients, right-hand sides,
and finite variable bounds are rational; finite binary floating-point
inputs mean their exact rational values. Each supplied affine row contains
only $y$ and at most one leaf. Rows containing only $y$ are allowed;
rows involving two different leaves are outside this class. Every
quadratic coefficient can have either sign. The same structure must hold
for the scalar support objective obtained from a proposed direction.

For fixed $y$, each leaf is constrained to an interval
\begin{equation}\label{eq:leaf-bounds}
L_i(y)\le x_i\le U_i(y),\qquad
L_i(y)=\max_{r\in\mathcal L_i}(\alpha_{ir}+\beta_{ir}y),\quad
U_i(y)=\min_{r\in\mathcal U_i}(\gamma_{ir}+\delta_{ir}y).
\end{equation}
The original bounds provide at least one function in each envelope.
Projecting each bounded center--leaf polygon onto the center gives a
closed rational interval. Their intersection, together with the center
bounds and center-only rows, is exactly the feasible center interval
$I$: conditional leaf feasibility is independent across leaves. An
empty intersection proves this support domain empty.

\begin{theorem}[Exact rational star support]\label{thm:star}
For the nonempty domain just described, the minimum of
\eqref{eq:star} is rational and has a rational minimizer. It is obtained
by a finite rational partition of $I$ on which each leaf has an affine
minimizer and the conditional optimum is quadratic in $y$. Both the
partition and the attained minimum can be reconstructed exactly from
the supplied rows, bounds, and coefficients.
\end{theorem}
\begin{proof}
Partition $I$ at all intersections of affine functions in the envelopes
of each leaf. On every open resulting interval, the active functions
$L_i$ and $U_i$ are fixed and affine.

If $d_i>0$, the unconstrained leaf minimizer is the affine function
$v_i(y)=-(e_iy+f_i)/(2d_i)$. The constrained minimizer is
$\max\{L_i(y),\min\{v_i(y),U_i(y)\}\}$. Further partition at
intersections of $v_i$ with its two active bounds. Each remaining
interval has one affine minimizing rule.

If $d_i\le0$, a minimizing leaf value is an endpoint. The difference
between the upper and lower endpoint objectives factors as
\[
q_i(y,U_i(y))-q_i(y,L_i(y))
=(U_i(y)-L_i(y))
\{d_i(U_i(y)+L_i(y))+e_iy+f_i\}.
\]
The first factor is nonnegative on the feasible center interval. The
second is affine on the current interval, so at most one additional
rational root determines the preferred endpoint; an identically zero
factor permits either endpoint. This also gives an affine minimizing
rule after refinement.

Take the union of all leaf breakpoints. Substituting the selected affine
rules into~\eqref{eq:star} gives a rational quadratic on every closed
piece. The rule remains feasible and minimizing at its endpoints:
bound switches coincide there, and endpoint-choice switches have equal
values there. Minimize each quadratic at its endpoints and, for a
positive square coefficient, its feasible stationary point. These
candidates are rational and exhaustive. The selected center and leaf
rules supply an attained rational minimizer. A singleton center interval
uses the same conditional scalar minimization at that point.
\end{proof}

If leaf $i$ has $m_i$ envelope lines, their pair intersections produce
at most $O(m_i^2)$ center breakpoints; a constant number of further
switches per envelope interval gives the same order. The union of
breakpoints has polynomial size. The present implementation deliberately
uses straightforward pair enumeration and scanning. Its arithmetic
cost is polynomial in the supplied row counts and number of leaves;
no claim of an optimal complexity bound is needed for this small-block
oracle. Unlike general elimination on a deeper graph, this construction
does not repeatedly nest piecewise messages along a tree.

The inherited implementation
\pathref{../../research-20261002-convexification/theory/quadratic_star.py}{theory/quadratic_star.py}
records the complete center partition, affine leaf rules, quadratic
value functions, and attained minima. Unsupported mixed terms or rows
raise an explicit error instead of being dropped. Replay rebuilds the
same exact object from trusted input. As with the polytope oracle, this
is exact support for a supplied direction; completeness of the
integration's direction search is a separate question.

\begin{corollary}[Sharp gain for specified pair directions]\label{cor:gain}
Let compact rational center--leaf polygons $P_i$ assemble into a nonempty
star, and fix rational quadratic pair objectives $q_i(y,x_i)$. Define
\[
\beta_i=\min_{P_i}q_i,\qquad
\beta_* = \min_{(y,x_i)\in P_i\ \forall i}\sum_iq_i(y,x_i),
\qquad \Delta=\beta_*-\sum_i\beta_i.
\]
Then $\Delta\ge0$ is the largest possible constant improvement in
the sum of these fixed pair inequalities. It is exactly computable by
the pair and star support algorithms. Moreover, $\Delta>0$ if and
only if the center projections of the pair minimizing sets have no
common point.
\end{corollary}
\begin{proof}
Every pair objective is at least its individual minimum, proving
$\Delta\ge0$. Equality holds precisely when one feasible assembled
point minimizes every pair objective. Such a point exists precisely
when their minimizing center sets intersect, since leaves can then
be chosen independently. If no such point exists, compactness makes
the attained assembled minimum strictly larger. The attained value
$\beta_*$ is the largest valid constant in this summed direction.
\end{proof}

This gives a precise diagnostic for a proposed merger of directions.
It does not choose the normals or predict whether computing the gain
will improve total runtime. In the example below, both pair minima
are zero and the exact gain is $1/128$.

The box-only case is a specialization of Del Pia and Khajavirad's
forest-QP algorithm, which gives an $O(n^2)$ arithmetic bound for general
forests \cite{DelPiaKhajavirad2026}. Piecewise-affine responses of
strictly convex parametric quadratic programs are also classical
\cite{BemporadEtAl2002}. The additional domain handled explicitly here
consists of center--leaf affine rows with all signs of scalar leaf
curvature. That scope distinction and its exact implementation are
documented; a first-result claim does not follow from a bounded source
search.

% Adapted from the 2026-10-02 report; inherited result and implementation.
\section{What overlapping pair hulls miss}\label{sec:overlap}

Exact pair support does not make intersected pair hulls the exact global
hull. The failure already occurs on a path $x$--$y$--$z$ in the unit
cube, even if every square and edge product is present. Define
\[
H_{xy}=\conv\{(x,y,x^2,xy,y^2):0\le x,y\le1\}
\]
and $H_{yz}$ analogously. Let $R$ intersect these two hulls while
identifying the shared first and second moments of $y$. The full sparse
quadratic hull is
\[
H=\conv\{(x,y,z,x^2,y^2,z^2,xy,yz):0\le x,y,z\le1\}.
\]

\begin{proposition}[A strict pair-hull gap]\label{prop:overlap}
The point
\begin{equation}\label{eq:moment-point}
\bar v=(1/2,1/2,4/5,1/2,5/16,4/5,3/8,1/2)
\end{equation}
belongs to $R$ but not $H$. The quadratic function
\begin{equation}\label{eq:D}
D=(y-1/4-x/2)^2+(y-5z/8)^2+x(1-x)+z(1-z)
\end{equation}
has minimum $1/128$ on the cube, whereas its lifted minimum over $R$
is zero.
\end{proposition}
\begin{proof}
The left pair measure assigns probability $1/2$ to each of
$(x,y)=(0,1/4),(1,3/4)$. The right pair measure assigns probabilities
$1/5,4/5$ to $(y,z)=(0,0),(5/8,1)$. Both give
$\E y=1/2$ and $\E y^2=5/16$, and their remaining moments yield
\eqref{eq:moment-point}. Thus the point belongs to $R$.

In a global representing measure, $\E[x(1-x)]=\E[z(1-z)]=0$
would force $x,z$ to be binary almost surely. The given moments also
force $\E(y-1/4-x/2)^2=\E(y-5z/8)^2=0$. The first identity requires
$y\in\{1/4,3/4\}$, and the second $y\in\{0,5/8\}$, almost surely.
Their supports are disjoint, so no global measure exists.

For fixed $y$, the square coefficients of $x$ and $z$ in $D$ are
$-3/4$ and $-39/64$. A minimum therefore has binary leaves. At those
endpoints the last two terms vanish. Set $a=1/4+x/2$ and $b=5z/8$.
Then
\[
(y-a)^2+(y-b)^2
=2\bigl(y-(a+b)/2\bigr)^2+(a-b)^2/2.
\]
The midpoint lies in $[0,1]$. The smallest difference between
$\{1/4,3/4\}$ and $\{0,5/8\}$ is $1/8$, giving $1/128$,
attained at $(x,y,z)=(1,11/16,1)$. Each pair summand of $D$ is
nonnegative on its pair box, and both have zero expectation under the
local witnesses. Their lifted minimum over $R$ is therefore zero.
\end{proof}

Expanding $D\ge1/128$ gives a cut involving only the existing sparse
quadratic coordinates. With $m_u,s_u,c_{uv}$ denoting the first,
square, and product coordinates, it is
\begin{equation}\label{eq:overlap-cut}
2s_y-\tfrac12m_y-c_{xy}-\tfrac54c_{yz}
+\tfrac54m_x-\tfrac34s_x+m_z-\tfrac{39}{64}s_z
\ge-\tfrac7{128}.
\end{equation}
It excludes~\eqref{eq:moment-point} by $1/128$, with no $xz$ product.
The star oracle of Theorem~\ref{thm:star} computes this support value
exactly. Scaling the objective scales the gap; $1/128$ is an exact
diagnostic value, not a universal limit on missing strength.

This comparison is against exact \emph{pair} hulls, not against a
fully coupled SDP relaxation. A dense moment matrix together with
the missing nonedge product and its McCormick inequalities excludes
this witness too. At the proposed moments,
$\operatorname{Var}(y)=1/16$, $\operatorname{Var}(x)=1/4$,
$\operatorname{Var}(z)=4/25$, and the two edge covariances saturate
the positive Cauchy--Schwarz bounds. A global PSD moment matrix would
therefore require $\operatorname{Cov}(x,z)=1/5$, hence
$\E[xz]=3/5>\E[x]=1/2$, contrary to $xz\le x$. The star cut
repairs this example without adding that nonedge coordinate; general
dominance over SDP or RLT is not asserted.

There is a simple family behind this example. Suppose
$a,a+b,c,c+d\in[\ell,u]$ and let
$\delta=\min\{|s-t|:s\in\{a,a+b\},t\in\{c,c+d\}\}$.
For $0\le x,z\le1$, $\ell\le y\le u$,
\begin{equation}\label{eq:overlap-family}
(y-a-bx)^2+(y-c-dz)^2+b^2x(1-x)+d^2z(1-z)
\ge\delta^2/2.
\end{equation}
The left side is affine in each leaf after expansion, so its minimum
occurs at binary leaves. The preceding midpoint calculation proves
both validity and sharpness. Automatic selection of the four parameters
is not supplied merely by this family; the star support algorithm is
the implemented general mechanism within its stated class.

\subsection{The stronger interface that makes tree gluing valid}

The obstruction concerns inconsistent representing distributions, not
an error in the individual hulls. Under a running-intersection tree of
compact block domains, local probability measures glue if they have
the \emph{same complete marginal distribution} on every adjacent
separator \cite{Lasserre2006}. Root the tree; retain the root measure
and attach each child's new variables using its conditional distribution
given the parent separator. Matching marginals preserve that child's
distribution, and running intersection prevents inconsistent assignments
of previously introduced variables. Induction constructs a global
representing measure. Conditional distributions exist for these compact
Euclidean Borel spaces; finite-support gluing is ordinary conditional
probability and finite summation.

For a binary scalar separator, its mean determines the complete
distribution. For a separator known to lie in a fixed set of $k$
distinct scalar values, matching moments of orders $0,\ldots,k-1$
determines atom probabilities through an invertible Vandermonde system.
Neither observation applies to unrestricted continuous separators with
only first and second moments. Means of several binary separator
variables also do not determine their joint distribution.

The counterexample does not contradict exactness of the McCormick
relaxation for a \emph{pure bilinear} graph on a forest. On the unit
box, an edge with means $m_i,m_j$ and product $w_{ij}$ has binary
probabilities
\[
p_{11}=w_{ij},\quad p_{10}=m_i-w_{ij},\quad
p_{01}=m_j-w_{ij},\quad p_{00}=1-m_i-m_j+w_{ij}.
\]
The McCormick inequalities are their nonnegativity conditions. Adjacent
edges share a binary marginal fixed by the common mean, so these
measures glue along the forest. Prescribed square coordinates can
forbid this binary choice, as they do in~\eqref{eq:moment-point}.

The companion
\pathref{../../research-20261002-convexification/theory/overlap-review.md}{theory/overlap-review.md}
also gives a nonlinear constrained example that survives separate
block interval propagation and is repaired by one additional separator
statistic. That example illustrates a different possible extension;
it is not claimed to defeat global presolve or to be an automatic
statistic-selection algorithm. The inherited implementation supplies the concrete alternative of
selected larger star support cuts. The current exact polytope oracle also
handles their union whenever it fits its explicit enumeration budget.

\section{Complete separation to a positive tolerance}\label{sec:separation}

Exact support for a chosen direction does not ensure that a bounded
numerical search finds the direction needed to separate a query. This
section supplies a finite fallback with a complete answer contract. Its
potentially large cost is part of the result.

Let $P$ be a nonempty compact rational polytope, let
$F:P\to\R^k$ be a vector of rational polynomials, let
$q\in\Q^k$ be a query, and choose a positive rational tolerance
$\epsilon$. A designated subset $J$ of coordinates may have nonnegative
support coefficients. Define
\[
C_J=\{c\in[-1,1]^k:c_j\ge0\ (j\in J)\},\qquad
U_J=\conv F(P)+\operatorname{cone}\{e_j:j\in J\}.
\]
For $J=\varnothing$, this is ordinary graph-hull separation, provided
the coordinate features of $x$ are included in $F$ when required.
For nonempty $J$, the cone expresses inequality-side slack, as in
Theorem~\ref{thm:closure}. Set
\[
\psi(c)=\min_{x\in P}c^\top(F(x)-q).
\]

\begin{lemma}[Distance and direction continuity]\label{lem:distance}
The distance in feature coordinates satisfies
\[
\operatorname{dist}_1(q,U_J)=\max_{c\in C_J}\psi(c).
\]
If rational bounds $\ell_j\le F_j(x)\le u_j$ hold on $P$, put
$R=\sum_j\max\{|\ell_j-q_j|,|u_j-q_j|\}$. Then
$|\psi(c)-\psi(d)|\le R\norm{c-d}_\infty$ on $C_J$.
\end{lemma}
\begin{proof}
For any $z\in U_J$, any $c\in C_J$ satisfies
$\psi(c)\le c^\top(z-q)\le\norm{z-q}_1$, giving one inequality.
The set $U_J$ is nonempty, closed, and convex by the compact-plus-cone
argument in Theorem~\ref{thm:closure}. If its distance from $q$ is
positive, separate it from the open $L_1$ ball about $q$ with that
radius. Scale the separating functional to infinity norm one. Its
finite lower support requires $c_j\ge0$ on cone coordinates, and
the support value of the ball yields the reverse inequality. At distance
zero, $c=0$ and the first inequality give equality.

For the continuity bound, every $x\in P$ satisfies
$|(c-d)^\top(F(x)-q)|\le R\norm{c-d}_\infty$.
Taking minima preserves the same uniform error in both directions.
\end{proof}

The field's support/separation equivalence is classical. The familiar
polynomial-time weak-oracle equivalence requires a specified encoded
convex-body model, including its inner and outer radius information
\cite{GrotschelLovaszSchrijver1981}. The elementary construction here
uses no inner ball and works for lower-dimensional hulls. Its enumeration
cost is not a replacement polynomial-time theorem.

\begin{theorem}[Finite positive-tolerance answer]\label{thm:finite-separation}
Suppose each chosen normal $c\in C_J$ can be assigned a feasible
$p_c\in P$ and a certified lower support bound $L_c$ satisfying
\[
L_c\le\min_{x\in P}c^\top F(x)\le c^\top F(p_c)
\le L_c+\delta,
\qquad 0\le\delta<\epsilon.
\]
Choose an integer
$N\ge\max\{1,\lceil2R/(\epsilon-\delta)\rceil\}$.
Enumerate the Cartesian normal grid with coordinates
$-1+2i/N$ for unrestricted components and $i/N$ for components in $J$,
where $i=0,\ldots,N$. If any $L_c>c^\top q$, return the valid rational
cut $c^\top z\ge L_c$. Otherwise
$\operatorname{dist}_1(q,U_J)\le\epsilon$.
\end{theorem}
\begin{proof}
A returned cut is valid by its lower support bound and nonnegative cone
coefficients, and it strictly excludes $q$. If none is returned, every
grid normal satisfies
$\psi(c)\le c^\top(F(p_c)-q)\le\delta$.
Every normal in $C_J$ is within infinity distance at most $2/N$ of
a grid point. Lemma~\ref{lem:distance} bounds its support gap by
$\delta+2R/N\le\epsilon$, and the distance identity proves the claim.
\end{proof}

For quadratic features, Theorem~\ref{thm:polytope} supplies exact support
and a rational minimizer, so $\delta=0$. A general rational polynomial
graph also admits a finite support approximation. Let
$v_1,\ldots,v_s$ be all vertices of $P$ and let
\[
D_\infty=\max_i\left(\max_{x\in P}x_i-\min_{x\in P}x_i\right),
\qquad
L\ge\sum_{j=1}^k\sum_{i=1}^d
\max_{x\text{ in supplied box}}|\partial_iF_j(x)|.
\]
Rational interval polynomial evaluation supplies a finite $L$.
For a positive integer $M$, enumerate
\[
S_M=\left\{\sum_{i=1}^s\frac{n_i}{M}v_i:
n_i\in\mathbb Z_+,\ \sum_i n_i=M\right\}.
\]
All its points are feasible, even on an affine equality domain of zero
ambient volume. This avoids the failure of filtering a rectangular
floating-point grid for exact equality feasibility.

\begin{lemma}[Rational domain net]\label{lem:domain-net}
Every $x\in P$ is within infinity distance
$D_\infty(s-1)/M$ of a point of $S_M$. Consequently, for every
$c\in C_J$,
\[
\min_{x\in S_M}c^\top F(x)-\delta_M
\le\min_{x\in P}c^\top F(x)
\le\min_{x\in S_M}c^\top F(x),
\qquad \delta_M=LD_\infty(s-1)/M.
\]
\end{lemma}
\begin{proof}
Express $x=\sum_i\alpha_iv_i$ as a convex combination. For the first
$s-1$ coefficients put $n_i=\lfloor M\alpha_i\rfloor$ and put
the remaining mass in $n_s$. The difference from the resulting sample
is $\sum_{i<s}(\alpha_i-n_i/M)(v_i-v_s)$, whose infinity norm is
at most $D_\infty(s-1)/M$. The derivative bound implies that
$c^\top F$ is $L$-Lipschitz for the infinity norm on the box.
Apply that inequality to a minimizer and its sample. The reverse bound
uses $S_M\subset P$.
\end{proof}

Choose $M\ge\max\{1,\lceil2LD_\infty(s-1)/\epsilon\rceil\}$.
Then $\delta_M\le\epsilon/2$ and Theorem~\ref{thm:finite-separation}
applies. A singleton or constant graph causes no division by zero.
The domain enumeration has at most $\binom{M+s-1}{s-1}$ candidates;
the normal grid has $(N+1)^k$ directions. The vertex enumeration and
each exact support calculation have their own costs. These bounds show
finite termination and explicitly expose the dependence on dimension,
vertex count, coefficient size, and tolerance. They are generally too
large for a solver callback.

The implemented API
\pathref{../solver/separation.py}{solver/separation.py}
uses numerical proposals and exact convex-combination checks first.
Successful proposals may avoid either grid. A rational convex combination
$z_*\in\conv F(P)$ proves the ordinary near-hull result when
$\norm{z_*-q}_1\le\epsilon$; with cone coordinates the check is
\[
\sum_{j\notin J}|z_{*j}-q_j|+
\sum_{j\in J}\max\{z_{*j}-q_j,0\}\le\epsilon.
\]
Numerical weights are clipped to be nonnegative and renormalized in
exact arithmetic, and every graph sample is checked against the exact
domain. An inaccurate LP can therefore fail to help but cannot itself
authorize an incorrect geometric verdict.

There are four mathematical return statuses: \texttt{cut},
\texttt{within\_tolerance}, \texttt{empty\_domain}, and
\texttt{unresolved}. The first carries an exact rational separating row;
the second a distance certificate; the third an exact empty-polytope
calculation. The fourth has no geometric implication. In complete mode,
the enumeration budgets are removed and the rational algorithm terminates
for every admitted finite input. In bounded mode, an exhausted budget
returns \texttt{unresolved}. Unsupported or malformed inputs are explicit
input errors, not one of the successful mathematical statuses.
The bounded API limits support-call counts and polynomial domain-net
candidates. It does not itself impose a wall-clock limit or cap the
initial vertex enumeration and each exact quadratic face enumeration.
The deployed callback uses a separate support wrapper with explicit
face-count caps.

A near-hull grid certificate needs only one feasible point per direction
whose checked gap is at most $\delta$, together with verified grid coverage
and $\delta+2R/N\le\epsilon$. Replay need not trust the direction LP or
its claim to have optimized those points. A cut, in contrast, needs a
lower bound on the whole domain. These distinct evidence requirements
are checked separately. An optional binary64 version is corrected using
feature bounds; its certificate records whether rounding preserved strict
separation. Exact separation can survive when no useful numerical row does.

For direct-row separation use $F=(x,g(x))$, the query
$(\bar x,b-L\bar z)$, and nonnegative normal coordinates for $g$.
An exact returned row maps through Proposition~\ref{prop:aggregate}.
A \texttt{within\_tolerance} result measures distance in these feature/slack
coordinates to $U_J$; it is not automatically the same distance in the
original model variables to the projected feasible set. The practical
solver loop uses its explicitly bounded proposal policy. The complete
API establishes an available finite completion of this specified
separation problem, not an assertion that every callback runs that fallback.

\section{The implemented model, cut, and replay contract}\label{sec:implementation}

The common model builder and the cut generator have different responsibilities.
The builder preserves the admitted original model at the submitted expression
boundary. A cut can then use only a conservative domain justified from that
model. Declining a difficult block leaves the original native model available;
it does not authorize changing the function, discarding part of a row, or
declaring the point feasible for the nonlinear constraints.

\subsection{Source arithmetic and native expression binding}

Every loaded finite binary64 number is interpreted as its exact rational
value. Ordinary polynomial assembly can nevertheless change intermediate
values: the source coefficient in $(0.1x)^2$ is the exact square of the
binary64 leaf, whereas multiplying its coefficient in binary64 rounds that
square again. Likewise a native constraint constructor can move a constant
into its sides with a rounded subtraction. Neither difference is accepted
merely because it is small.

The builder first assembles the ordinary native expression, then compares
its exact interpreted value with the source expression by symbolic subtraction
and expansion. If they differ, it constructs an explicit generic expression
DAG with separate constant children and checks again. This retains the
original coefficient operations instead of folding them into changed
coefficients. A remaining difference gives an explicit model-admission
refusal. Ranged constraint sides and objective sense and constants are part
of the same check. The comparison has no numerical tolerance.

This boundary is the submitted PySCIPOpt expression DAG interpreted over
the reals with its exact binary coefficients. SCIP's later simplification,
presolve, LP arithmetic, and feasibility tolerances are outside this symbolic
binding claim. A successful model build is not a proof of the complete solve.

\subsection{Original affine implications and exact domains}

Original affine rows can prove finite or one-sided bounds that are absent
from the declared box. For a row $\sum_i a_ix_i\le b$, the exact
lower activity of all terms other than $j$ implies
\begin{equation}\label{eq:bound-propagation}
a_jx_j\le b-\sum_{i\ne j}\min\{a_i\ell_i,a_iu_i\}.
\end{equation}
Division by the nonzero $a_j$ gives the corresponding upper or lower
bound; a lower row side is handled after negation. Integer and binary
semantics permit exact integer rounding where justified. Every recorded
deduction identifies the original row, side, target variable, and previous
and new bounds. Replay rebuilds the premises and recomputes the sequence.
A budget stop leaves a valid prefix; it does not claim propagation closure.
When exported, lower bounds round outward downward and upper bounds upward.

\begin{proposition}[Affine-bound provenance]\label{prop:affine-bounds}
Starting from valid declared and variable-type bounds, every bound obtained
by a finite sequence of~\eqref{eq:bound-propagation} and valid integer
rounding contains all original feasible points.
\end{proposition}
\begin{proof}
At a feasible point, every other term is at least its lower activity
under valid prior bounds. Subtracting those bounds from the original row
gives~\eqref{eq:bound-propagation}. Dividing with the correct sign
preserves it; rounding to an integer endpoint preserves all integer
solutions. Induction proves the statement for the complete recorded prefix.
\end{proof}

The model builder also checks affine rows proportional to a source domain
argument. For example, the equality $x-y=1$ proves that $\log(x-y)$
has a positive argument even when both coordinates are unbounded. This
argument-level implication does not invent finite coordinate bounds.

Domain traversal visits every original source node before simplification.
A zero multiplier, cancelled quotient, or zero exponent cannot erase the
domain required by its child expression. One-sided exact intervals preserve
useful information: $x\ge0$ proves $1+x\ge1$ without an upper bound.
When these bounds and original affine arguments do not prove a required
domain, an exact existential constraint preserves it:
\[
\begin{array}{c|c}
\text{required domain}&\text{additional constraint}\\\hline
d\ne0&du=1,\quad u\in\R\\
d>0&du=1,\quad u\ge0\\
d\ge0&d\ge0.
\end{array}
\]
In each of the first two rows, choosing $u=1/d$ proves the forward
projection and the product equality proves the converse; with $u\ge0$
it implies $u>0$ and $d>0$. These are exact strict-domain descriptions,
with no arbitrary positive epsilon. Witness variables can be unbounded.
The argument expression is itself compiled and checked against the source.
Shared witnesses are permitted only for the same argument and requirement,
and every nested original occurrence is still visited.

All solver modes use the same guards. They preserve domains such as those
hidden in $0\log x$, $x/x$, $(\sqrt{x})^2$, or $(\log x)^0$.
They can also change numerical difficulty and construction time; those
costs are measured. A strict source domain can give an unattained infimum.
Numerical optimality status and finite feasibility tolerances do not prove
attainment on that open domain.

The admitted variable-power rule is
$b(x)^{e(x)}=\exp(e(x)\log b(x))$ after positivity of $b(x)$ is
proved from original bounds and affine rows. The transformation is local
to that source node; a value such as $\log2$ remains a logarithm node,
not an approximated coefficient. If positivity is unproved, the builder
returns \texttt{unsupported\_variable\_power\_domain}. It does not
add a positive-base guard indiscriminately: doing so could remove valid
zero-base or negative-base points under broader exponent semantics.
Fixed exponents must fit the native power-node representation exactly.

The supported grammar and metadata are specified in
\pathref{../numerics/model-contract.md}{numerics/model-contract.md}.
Its separate admission artifact checks the seven previously refused
application models before optimization. Admission demonstrates broadened
applicability of this importer; performance and original-incumbent residuals
are different evidence.

\subsection{Block discovery and bounded practical separation}

The current integration groups complete nonlinear remainders of original
row sides. Exact additive affine terms are separated and retained in the
elimination formula. Overlapping remainders can form a larger group up to
the admitted block size; a nonlinear remainder involving unsupported
variables is declined as a whole. Original affine sides supported on the
chosen variables supply additional domain rows. Dropping a row from the
support domain enlarges it and can weaken a bound; extraction must not
change or drop terms from the nonlinear remainder multiplied by that
row multiplier.

Discovery is lazy, so a model solved in presolve does not pay block-search
cost. A numerical sample LP proposes a normal for the linear coordinates
and nonnegative multipliers for consistently signed nonlinear sides. Exact
support is then computed for the actual direction. Where available, its
rational minimizer becomes a new sample for a bounded exchange. Sampling
and numerical LP output choose work; they establish no lower bound by
themselves. Exact polytope vertices and their centroid seed proposals on
thin affine domains. Their conversion to floating-point samples is still
only a numerical proposal, not a proof that the converted points remain
exactly feasible.

The practical support wrapper admits general quadratic blocks up to
dimension four under a complete-subset budget, then uses inherited supported
polynomial, elementary, or constrained-star kernels where applicable.
The default standalone wrapper allows 5,000 subsets. The solver callback
sets a limit of 2,000 subsets per exact polytope calculation; an unmet
limit yields no cut from that calculation.
Supported higher-degree polynomial and elementary bounds remain sound
without an exact support claim. Partial-domain features whose complete
support domain cannot be safely covered are declined. If algebraic
cancellation yields a polynomial extension, support on a larger box is
valid for original feasible points, but a closeness certificate for that
extension is not a membership certificate for the original partial graph.

Three solver modes use the common original-model builder:
\begin{itemize}
\item \texttt{baseline} supplies no additional aggregation separator;
\item \texttt{all} uses the bounded search on all admitted blocks;
\item \texttt{auto} uses a fixed structural admission rule and requires a
larger final normalized violation before inserting a cut.
\end{itemize}
The mode name \texttt{all} means all admitted blocks subject to its caps,
not all mathematical support directions. The exact finite completion API
of Section~\ref{sec:separation} is available separately; callback time
limits are not relabeled as a complete separation result.

The frozen \texttt{auto} rule admits quadratic blocks containing a signed
quadratic not proved convex, and either a non-axis affine domain row or
at least two distinct nonlinear source rows. Two support failures deactivate
a block for the remaining callbacks. This is a heuristic spending decision,
not a certificate that the block cannot help at a later point. Both cut
modes permit at most 32 blocks, six nonlinear sides and 16 affine domain
sides per block, three root callbacks, 12 total cuts, four cuts per
callback, 24 support calls, and three LP exchange attempts per block
and callback. Callback work receives the smaller of one second and five
percent of the requested total budget, checked between operations.
Nonpreemptive work can overrun that allowance; the campaign's hard process
limit is a separate safeguard. The final violation divided by
$\max\{1,\sum_j|\widehat c_j|\}$ must exceed $10^{-5}$ in
\texttt{all} and $5\cdot10^{-4}$ in \texttt{auto}.
The corrected discovery routine checks the existing deadline between
source rows, additive terms, and group steps. An expired allowance
discards unfinished discovery and records
\texttt{discovery\_incomplete} and \texttt{budget\_exhausted}; it
produces neither a coverage claim nor a hull-membership conclusion.
Terms are converted to polynomials using only their own symbols.
Nonrational coefficients and oversized terms that cannot be converted
remain in the nonlinear remainder, so an unsupported remainder is
declined without changing the native row. The prospective experiment
and the separate correction validation are distinguished below.

Each accepted cut follows the same sequence: prove support for the actual
direction; reconstruct source-side elimination exactly; combine repeated
original-variable coefficients; apply the bound correction of
Proposition~\ref{prop:elimination-round}; inspect the actual stored row;
and submit it as a global root cut. All domain rows and rounding bounds
must be original global implications. No incumbent cutoff or local
presolve bound is used as an unmarked global premise.
The row-conversion implementation is conservative on unbounded coordinates:
although Proposition~\ref{prop:elimination-round} needs only the bound
in the error's direction, the frozen code requires both finite bounds
whenever that coordinate has nonzero rounding error. Exact coefficients
need no bounds. This narrower admission can reject a mathematically safe
row; it cannot strengthen an unsafe one.

There is no nonlinear feature-equality reformulation that differs across
modes. The common objective epigraph and any original-domain witnesses
are recorded in every mode. The new comparison therefore needs no
separate duplicate reformulation control. The earlier four-mode experiment
retains its control because its formulation did change.

\subsection{Replay and the trust base}

A cut record binds the original serialized instance, signed source row
identities, variable identifiers, source nonlinear expressions, proved
affine bounds, chosen support domain, final direction, exact support
certificate, source-side elimination, rounding correction, and actual
binary64 row. Fresh-process replay must derive these expected objects
from the archived original model. A stored explanation or producer/checker
agreement without that binding would be insufficient.

Exact arithmetic and interval routines are shared between generation and
replay. The checker is independent of the numerical direction LP and
solver search, but is not a second formally verified arithmetic
implementation. Separate analytic examples, alternative-oracle checks,
domain adversaries, and deliberate certificate corruptions provide distinct
implementation evidence. They complement the proofs rather than replace
them. Numerical original-coordinate residual checks and reference-bound
comparisons remain explicitly separate from cut validity.

\section{Prospective evaluation and verification}\label{sec:evidence}

\subsection{The earlier result is retained}

The first experiment completed 316 scheduled records. Its 24 selected
application models included 20 admitted models; baseline and matched
reformulation control solved 19, while both cut modes solved 18.
Automatic selection reduced direction LPs from 1,073 to 323 and callback
time from 5.59 to 3.05 seconds relative to unconditional activation,
but did not recover the lost solve. On 13 synthetic cases, the control
and cut modes solved all 13 while baseline solved 12, so the extra solve
could already be explained by reformulation.

All 1,082 recorded added cuts passed the first campaign's replay and all
270 returned incumbents passed its numerical residual checks. Eight
construction-error records lacked model/cut logs; those missing logs were
excluded from successful replay coverage and retained as unknown counts.
These historical results justify investigating a different integration,
not a favorable claim about its performance in advance.

\subsection{Frozen new comparison}

The new
\pathref{../experiments/protocol.md}{prospective protocol}
was fixed before primary outcomes. It excludes the prior convexification
campaign names and selects 30 cached MINLPLib OSiL models, ten each from
convex models, nonconvex continuous models, and nonconvex models with
integer variables. Eligibility uses size limits and syntax parsing, not
success of the new importer or cuts. Within each stratum, a fixed
SHA-256 ordering selects models. The 422 eligible models have stratum
counts 85, 208, and 129. ``New holdout'' refers to the named earlier
convexification campaigns, not to every use of these public instances
elsewhere in the repository.

Every selected failure remains in its denominator. The common builder
is used in \texttt{baseline}, \texttt{all}, and \texttt{auto}. Primary
full runs use seed zero, one SCIP thread, one BLAS/OpenMP thread, a
30-second soft total integration budget, and relative gap tolerance
$10^{-4}$. Model loading, analysis, building, cut generation, and solving
are included in the integration budget; process imports and independent
result checking are counted in outer wall time. The hard process limit
for these runs is 45 seconds. Nonpreemptive calls can overshoot a soft
budget and must be reported.

The fixed schedule comprises 90 new-holdout full runs, 30 historical
diagnostic runs, 39 full mechanism runs, 105 root-only runs, and 18
seed-one repeats: 282 jobs. The 13 mechanisms use ten-second soft
budgets; root-only runs use five seconds and one node. Mode order rotates
by model index. One worker runs at a time on a host shared with unrelated
jobs, so small timing differences remain descriptive. The 150-minute
aggregate emergency cap exceeds the sum of all scheduled hard process
timeouts; any unstarted jobs would still be recorded.

Frozen source copies and hashes, exact tagged parsed models, original
OSiL files, mode configurations, solver status, numerical bounds, nodes,
original-variable incumbents, work statistics, and cut witnesses are
retained. Every returned incumbent is evaluated against original rows,
bounds, domains, integrality, and objective by a separate evaluator at a
scaled tolerance of $10^{-5}$. Recorded root and final bounds are compared
with known synthetic optima or archived feasible MINLPLib objectives.
Such a conflict is a diagnostic flag, not by itself a proof about which
numerical calculation is wrong. Fresh-process replay checks each saved
cut against its original model and actual stored row, with tampering
controls and explicit accounting for missing logs.

The campaign source snapshot is distinct from the final working tree.
The first post-freeze changes rejected cut sides at SCIP's infinity
sentinel, checked that side in stored-row auditing, and caught failed
heuristic exchange-sample evaluations. The campaign then exposed four
worker errors in cut modes on \texttt{chp\_partload} and
\texttt{waterno2\_06}: affine splitting passed an unnecessarily large
list of symbolic generators to polynomial conversion. Discovery on a
large case also overran its allowance without checking between rows.
The repair uses each term's actual generators, safely declines terms
whose rational polynomial conversion is unavailable, and checks the
existing work allowance between discovery operations. These repairs have
their own targeted review and matched validation runs.

The archived prospective experiment keeps its source and every original
result. Its timings describe that frozen implementation. Missing worker
cut logs remain unknown. A separate supplement repeats the affected
cases with all comparison modes after the correction; it does not replace
the original outcomes or turn outcome-selected repairs into a new holdout.
The final reconciliation records exact repair selection, source identities,
work-budget effects, and replay coverage separately.

\subsection{Prospective campaign results}

The campaign completed all 282 scheduled jobs in 1,338.6 seconds of outer
wall time, with no hard process timeouts or unstarted jobs. There were
106 optimal statuses, 74 gap-limit statuses, 31 time limits, 67 node
limits, and the four retained worker errors described above. All 30 new
holdout models were admitted in every mode. The same 25 were numerically
solved in every mode; the same five remained unsolved. Here solved means
an optimal or gap-limit status with an incumbent passing the independent
original-model numerical checks.

\begin{center}
\begin{tabular}{lrrrr}
\toprule
Mode & Solved/selected & Added cuts & Cases with cuts & Total seconds\\
\midrule
Native baseline & 25/30 & 0 & 0 & 170.57\\
All admitted blocks & 25/30 & 38 & 7 & 187.02\\
Automatic selection & 25/30 & 10 & 3 & 187.12\\
\bottomrule
\end{tabular}
\end{center}
Total seconds sum the recorded integration and preparation time; they
exclude interpreter imports and independent result checking. Shared-host
timing is descriptive. The cut modes establish no solved-count gain in
this population and take more summed time in these runs.

Automatic selection reduced candidate LPs from 186 to 66 and support
calls from 138 to 51. It did not substantially reduce callback time:
27.63 seconds for \texttt{all}, versus 27.00 for \texttt{auto}.
Discovery accounted for 24.41 and 24.63 seconds respectively, preceding
most of the admission policy's savings. Of 30 full runs per cut mode,
28 invoked discovery, 12 found no supported blocks, and two solved
without calling this separator. The admitted groups contained 156
possibly overlapping blocks, of which 63 satisfied the structural
automatic rule; 186 source sides were declined. These are observed
coverage and work counts, not unique variables or a causal decomposition
of native relaxation strength.

At the prespecified scaled dual-comparison tolerance $10^{-6}$, final
bounds for \texttt{all} were better on two models, tied on 21, and worse
on seven; \texttt{auto} had two better, 24 tied, and four worse. Both
improved bounds occurred on models already solved in all modes.
Among the five common unsolved cases, both cut modes had zero better,
one tied, and four worse bounds. None of those five received an added
cut. For example, \texttt{graphpart\_clique-40} finished with lower
bound 438 for baseline and 392 for each cut mode. Discovery and callback
cost, time limits, and shared-host variation prevent attributing every
difference to one cause; these results do not show useful added strength
on the common hard cases.

The root-only application comparison gave one better, 26 tied, and three
worse bounds for \texttt{all}, versus one better, 28 tied, and one worse
for \texttt{auto}. The seed-one repeats solved all six
selected models in every mode and had six tied final bounds in each
comparison. The complete root and repeat tables remain in
\pathref{../experiments/campaign-v2/summary.json}{the campaign summary}.
The repetitions are too small to establish a population-wide timing effect.

Every mode solved 12 of the 13 full synthetic cases. A positive local
mechanism is retained: on the root-only \texttt{simplex\_quadratic\_vector}
case, both cut modes added two cuts and reached an optimal status with
bound approximately $-0.50000002$, versus baseline's node-limit bound
$-0.50027374$. Its true optimum is $-1/2$. This is evidence that the
row-domain support can supply useful information on a constructed case;
it is not an extra full-solve success on the application sample.

The ten historical diagnostics gave five solved cases in each mode.
The four cut-mode failures occurred after model construction, so they
must not be reported as importer refusals. The successful admission
checks for the seven earlier refused models likewise establish importer
coverage, not a guarantee that every subsequent separator operation
finishes successfully. The repaired runs below address those errors
and the discovery-budget defect directly.

Fresh-process replay passed all 123 recorded cuts, including their actual
stored SCIP rows, and verified 156 frozen source/input hashes. All 14
tampering controls were rejected. The four error records have no complete
cut logs and retain unknown counts; the successful replay claim excludes
them. All 271 returned incumbents passed the independent numerical
original-model checks, with no recorded reference-bound conflicts. The
\pathref{../experiments/campaign-v2/replay.json}{replay record} and
independent experiment audit provide the exact coverage. These results
certify the recorded cut inequalities under the documented trust base,
not the numerical solver's final global bounds or complete search.

\subsection{Matched repair validation}

The correction cohort was fixed before any repaired optimization outcome.
It includes every original model/phase/seed group whose cut mode exceeded
its recorded discovery allowance or failed on the identified global-symbol
recursion path. The rule selected 25 groups and 75 matched jobs, including
all three modes for each group. The exact triggering values and original
record hashes are retained in
\pathref{../experiments/repair-plan.json}{the repair plan}; the
\pathref{../experiments/repair-protocol.md}{repair protocol}
keeps the original phase, seed, time and node limits, and mode order.
The corrected code has a separate source snapshot. No activation
threshold, instance-name filter, or selected population is retuned.

This is a regression cohort selected from observed failures and budget
violations, not another prospective holdout. Its results remain separate
and cannot be spliced into the 30-model table above. All 75 jobs completed
in 754.5 seconds, with no worker errors, missing cut logs, hard timeouts,
or unstarted jobs. The solved counts were identical across the three
modes within each cohort:

\begin{center}
\begin{tabular}{lrrrr}
\toprule
Matched cohort & Groups & Solved per mode & All cuts & Auto cuts\\
\midrule
Full application runs & 8 & 4 & 12 & 3\\
Full historical diagnostics & 7 & 4 & 10 & 0\\
Root application runs & 9 & 0 & 11 & 3\\
Seed-one repeat & 1 & 1 & 3 & 0\\
\bottomrule
\end{tabular}
\end{center}
Every group has a baseline, an \texttt{all}, and an \texttt{auto} run;
the table's solved column applies to each mode, not their sum. Baseline
adds no new cuts. There are no additional solves from either cut mode.
For the eight affected full application models, final bounds were zero
better, five tied, and three worse for \texttt{all}; \texttt{auto}
had zero better, six tied, and two worse. Summed integration times were
123.91, 124.86, and 124.82 seconds for baseline, \texttt{all}, and
\texttt{auto}. For the seven historical diagnostics both cut modes had
six tied bounds and one worse bound, on \texttt{waterno2\_06}.
All root and seed-repeat bound comparisons tied at the stated tolerance.

The four formerly crashing diagnostic runs now stop incomplete discovery
and continue native SCIP. Six discovery calls in the correction cohort
crossed their soft deadline by a single nonpreemptive unit of work;
all six reported incomplete discovery. The largest discovery excess
was 3.618 milliseconds. The largest full-callback excess was 29.529
milliseconds, and the largest total soft-budget excess was 23.611
milliseconds. These measurements support the checked stopping contract,
not a hard real-time bound or a claim that the discarded large blocks
were successfully analyzed.

All 42 recorded repair cuts passed fresh-process original-model and actual-row
replay, with 103 archived source/input hashes verified and all 14 tampering
controls rejected. All 67 returned incumbents passed the independent numerical
checks, with no reference-bound conflicts. The
\href{run:../experiments/repair-discovery-v1/results.md}{repair results},
\href{run:../experiments/repair-discovery-v1/replay.json}{replay}, and
\href{run:../reviews/experiment-review.md}{independent experiment review}
retain the original and repaired evidence separately. Across the two
current campaigns, all 165 recorded cuts passed replay; the four original
unknown logs remain outside that claim.

The operational recommendation remains native SCIP by default. The
corrected implementation resolves the observed execution and discovery-budget
defects, supplies optional certified cuts, and has no demonstrated general
solve advantage in these experiments. The exact support and complete
tolerance-separation APIs remain reusable mathematical components independent
of that negative activation result.

\subsection{Targeted implementation and document checks}

The initial combined topic-specific local test command passed 250 tests and
eight subtests. After the discovery repair, the focused integration, row
arithmetic, and replay command passed 81 tests and three subtests. The
\href{run:../VERIFICATION.md}{verification record}
gives the actual command and preserves an earlier test-collection import
failure and its correction. Separate independent support diagnostics
checked 40 analytic and adversarial polytope cases, including equality
domains, flat valleys, affine images of simplexes, and exhaustive Boolean
MaxCut controls. The separation review checked 21 cases covering finite-net
completion, missed interior normals, thin domains, coordinate cones,
polynomial support error, tiny rational gaps, and altered certificates.
The numerical direction LP is not an oracle for any of these checks.

The source/domain review includes cancelled restrictions, one-sided
unbounded intervals, original affine implications, positive-base variable
powers, generic-expression constant handling, and deliberate model-binding
corruption. Integration replay tests reconstruct actual SCIP rows from
original source sides and reject altered premises. Owner tests and internal
independent reviews cover different risks; agreement of producer and replay
alone would not establish the underlying theorem.

The report's \href{run:VERIFICATION.md}{document verification record}
gives its targeted LaTeX build and evidence-link checks. No project-wide
verification or CI inspection is part of these local records. Finite tests
support implementation confidence under the proved contracts; they do not
prove a whole numerical SCIP solve or a general performance advantage.

\section{Prior results and contribution boundary}\label{sec:literature}

Simultaneous convexification and support of vector graphs precede this
work. Tawarmalani develops inclusion certificates and simultaneous
convexification \cite{Tawarmalani2010}; Ballerstein's thesis treats
vector univariate graphs in a solver setting \cite{Ballerstein2013};
and Liers and coauthors develop simultaneous convexification for gas
networks \cite{LiersEtAl2021}. He and Tawarmalani's composite and
factorable relaxations already address multiple outer functions and
shared-variable graph constructions
\cite{HeTawarmalani2021,HeTawarmalani2022,HeTawarmalani2024}.
The recent axis-aligned framework also treats simultaneous factorable
graphs with feasible coupling \cite{ZhuHeTawarmalani2026}.
The present use of a weighted sum is an implementation basis, not a
new general convexification principle.

Low-dimensional quadratic hulls have exact conic and lifted
representations \cite{AnstreicherBurer2010}. Murty's exposition of total
face enumeration \cite[Section 2.9]{Murty1997} provides a direct
precedent for active-face stationarity. Rational quadratic certificates
are also classical; Vavasis proves
membership of rational quadratic programming in NP \cite{Vavasis1990}.
The new support implementation makes a direct finite enumeration and its
degenerate cases checkable. Box-constrained quadratic stars are already
inside the forest class of Del Pia and Khajavirad
\cite{DelPiaKhajavirad2026}; strictly convex parametric quadratic
programs have classical piecewise-affine optimizers
\cite{BemporadEtAl2002}. The inherited star construction states its
additional center--leaf affine-domain and arbitrary-curvature cases
explicitly. Neither its implementation nor a bounded source search proves
publication priority. Sparse SDP and moment approaches provide further
relevant exactness and overlap results
\cite{Khajavirad2026,Lasserre2006}.

The aggregation formula is weak Lagrangian duality
\cite[Sections 5.1--5.2]{BoydVandenberghe2004}. Quadratic aggregation
hull results have substantial additional structure and hypotheses
\cite{DeyMunozSerrano2022,BlekhermanDeySun2024}. Original-variable
Lagrangian rows also occur in optimization-based bound tightening
\cite{GleixnerEtAl2017}. Those rows can be redundant for the generating
LP yet useful for propagation. Here any additional strength must come
from joint optimization of the nonlinear aggregate over its domain.
Grouping small sets of quadratic terms before relaxation has established
solver precedents as well \cite{MisenerFloudas2012}.

Safe numerical cuts are established technology
\cite{CookEtAl2009SafeCuts}. The bound correction after rational-row
conversion in Proposition~\ref{prop:elimination-round} is the lower-row
form of the safe-approximation and aggregation principles of Eifler
and Gleixner \cite[Sections 2.2--2.3]{EiflerGleixner2024}.
Validated polynomial underestimation has a direct Bernstein precedent
\cite{GarloffSmith2008,GarloffJanssonSmith2003}. The current contribution
is the supported expression and domain contract, its exact or outward
arithmetic realization, the actual inserted-row binding, and the
reproducible record of use.

Complete positive-tolerance separation is also an established type of
problem. The classical weak optimization/separation equivalence
\cite{GrotschelLovaszSchrijver1981} is stronger algorithmically under
its stated encoded-body assumptions. Our finite normal and domain nets
give a direct proof and implementation without a full-dimensionality
assumption, at explicit and often impractical cost. The source audit
also checked the relevant 1984 corrigendum's scope
\cite{GrotschelLovaszSchrijver1984}; it addresses the later submodular
result, not the cited Theorem 3.1.

Native SCIP supplies substantial nonlinear relaxation and propagation
machinery \cite{BestuzhevaEtAl2025}. Its bivariate projected-domain
relaxations \cite{MullerSerranoGleixner2020}, quadratic handling, and
lifted cuts make a comparison against only scalar McCormick inequalities
too weak to establish practical value. The experiments therefore compare
complete native-model runs with the identical model builder and retain
unfavorable outcomes.

The primary-source audit is
\pathref{../literature/README.md}{literature/README.md}, with the exact
aggregation precedents in
\pathref{../literature/aggregation-prior.md}{literature/aggregation-prior.md}.
Its source manifest records versions, inspected passages, retrieval
limitations, and hashes. Inherited references not re-read in this
continuation retain their earlier provenance. For example, the current
Vavasis check covers the publisher abstract and metadata; the report's
finite candidate proof is given explicitly rather than attributed to an
unread proof. The research package makes no publication-priority claim.

\section{What completion does and does not imply}\label{sec:frontier}

The exact support and finite separation contracts resolve specified
mathematical problems. Their definitions also expose two limits that
cannot be removed by more careful cut implementation alone.

The documented development program is complete: proofs and exact APIs
cover the supported classes; admitted model semantics and exported rows
are checked; the finite separation fallback has defined success and
budget-stop meanings; and the prospective experiment plus matched repair
validation are retained and reviewed. The deployment decision is also
resolved by the evidence: keep native SCIP as the default and the extra
separator optional. Equal application solve counts and repaired execution
errors do not establish a general solver improvement.

\begin{proposition}[A general exact-support hardness boundary]\label{prop:hardness}
Exact support of arbitrary rational quadratic graphs on boxes, with
dimension as part of the input, is NP-hard.
\end{proposition}
\begin{proof}
Let an undirected graph have nonnegative integer weights $w_{ij}$ and put
\[
q(x)=-\sum_{\{i,j\}\in E}w_{ij}(x_i+x_j-2x_ix_j),
\qquad x\in[0,1]^{|V|}.
\]
The function is affine in each coordinate separately. Starting from a
global minimizer, move one coordinate at a time to a minimizing endpoint
while keeping the others fixed; its value cannot increase. A binary
global minimizer therefore exists. At a binary point,
$x_i+x_j-2x_ix_j$ is one precisely when the edge crosses the associated
cut. Thus $\min q=-\operatorname{MaxCut}(G,w)$. The construction is
polynomial in the encoded input, and weighted MAX CUT is NP-complete
\cite{Karp1972}.
\end{proof}

Consequently an unrestricted polynomial-time exact support method would
imply P=NP. The fixed-dimension and constrained-star support results do
not conflict with this boundary: one fixes dimension, and the other
restricts the interaction and constraint structure. The finite polynomial
graph separator is complete precisely because its potentially exponential
enumeration cost remains explicit. Hardness does not prohibit useful
heuristics, special classes, approximations, or better algorithms; it
prohibits treating unrestricted tractable support as a routine missing
step under the usual complexity assumption.

Second, complete separation of the chosen graph relaxation does not
complete the original feasible-set hull. The $x^2=1/4$ example in
Section~\ref{sec:aggregation} proves that distinction even in one
variable. Choosing stronger support domains or larger merged blocks can
help; the cost and exactness of those choices are new questions, not
consequences of finding every normal for the weaker relaxation.

There is also a practical distinction. A valid, complete, or stronger
oracle need not improve a solver's runtime. The native solver can already
derive comparable information, an expensive direction can arrive too
late, or the cuts can alter numerical behavior. Prospective comparisons
can finish with a negative deployment decision. That is evidence about
the investigated method, not a reason to omit its outcomes or rename
unsuccessful work as a success.

Finally, these certificates cover their bound model domains and submitted
rows. They do not verify all SCIP presolve substitutions, LP bases, node
bounds, spatial branching, integer inference, or original-model feasibility
and optimality in exact arithmetic. An end-to-end proof-producing MINLP
solver would need those additional components. Internal reviews and
replayed cuts establish the stated component-level evidence only.


\bibliographystyle{plain}
\bibliography{../literature/references}
\end{document}
